\documentclass[12pt]{article}
\usepackage[utf8]{inputenc}
\usepackage[T1]{fontenc}
\usepackage[brazil]{babel}
\usepackage[a4paper,margin=2.5cm]{geometry}
\usepackage{enumitem}
\setlist[itemize]{noitemsep,topsep=4pt}
\title{Sprint 3 - S3\_01\\Orientação a Objetos}
\author{Matheus Moreli Durão Silva\\Matrícula 12621CBS217\\Cibersegurança - Bacharelado - Noturno}
\date{2026}
\begin{document}
\maketitle

\section*{Entregável}
Hierarquia de classes, polimorfismo, JSON e recarga. Este documento registra o objetivo do card, a solução desenvolvida e a forma de conferência. Os arquivos de código e as evidências citadas permanecem na mesma pasta deste documento.

\section*{O que foi estudado e aplicado}
\begin{itemize}
\item Classe base Ativo e subclasses por tipo
\item Método sobrescrito para apresentar o tipo
\item Persistência e recarga dos objetos no JSON
\item Compatibilidade com o CRUD já existente
\end{itemize}

\section*{Como conferir}
O avaliador pode abrir os arquivos desta pasta, executar os exemplos quando houver código e comparar a saída com a explicação apresentada. Nos cards ligados ao Inventário de Ativos, a solução foi relacionada ao Trabalho 2 sem modificar requisitos já atendidos do CRUD.

\section*{Modelo e demonstração concreta}
\texttt{poo.py} define \texttt{Ativo} e as subclasses \texttt{Notebook}, \texttt{Servidor} e \texttt{Roteador}. Cada subclasse sobrescreve \texttt{exibir\_detalhes()} e chama \texttt{super()}. Percorrer uma lista mista executa o método de cada tipo sem seleção manual; \texttt{isinstance()} confirma a substituição por \texttt{Ativo}. \texttt{Vulnerabilidade} é objeto composto dentro do ativo, pois uma ocorrência não é um tipo de equipamento.
\section*{Persistência e prova de recarga}
\texttt{python poo.py} grava \texttt{objetos.json} como lista de objetos com campo \texttt{tipo}, lê o mesmo arquivo, recria as subclasses e usa assertivas para os IDs 1, 2, 3, os tipos e o vínculo da primeira vulnerabilidade. A saída real termina com ``JSON recarregado com IDs, subclasses e vínculos preservados''. O projeto principal tem seis subclasses e mantém a interface CRUD por funções de compatibilidade, com domínio em \texttt{Inventario}.
\section*{Conclusão}
A hierarquia, a composição e o ciclo JSON foram executados e podem ser repetidos pelo avaliador com a biblioteca padrão do Python.
\end{document}
